\documentclass{article}

% For arXiv preprint
\usepackage[preprint]{neurips_2026}

\usepackage[utf8]{inputenc}
\usepackage[T1]{fontenc}
\usepackage{hyperref}
\usepackage{url}
\usepackage{booktabs}
\usepackage{amsfonts}
\usepackage{amsmath}
\usepackage{amssymb}
\usepackage{nicefrac}
\usepackage{microtype}
\usepackage{xcolor}
\usepackage{graphicx}
\usepackage{algorithm}
\usepackage{algorithmic}
\usepackage{multirow}
\usepackage{subcaption}

\title{HEBBS: A Self-Tuning Memory Engine for AI Agents}

\author{%
  Parag Arora \\
  hebbs.ai \\
  \texttt{parag@hebbs.ai} \\
  \And
  Jasen Lam \\
  hebbs.ai \\
  \texttt{jasen@hebbs.ai} \\
}

\begin{document}

\maketitle

\begin{abstract}
AI agents that accumulate memories over weeks face two unsolved problems: memory degrades as contradictions compound and stale facts persist, and retrieval degrades because no single configuration is optimal across query types.

Existing agent memory systems address the first problem with better search over static stores. None address the second: agents retrieve the same way regardless of what they are retrieving, with no mechanism to evaluate or improve their own recall.

We present \textsc{Hebbs}, a memory engine where the agent controls how it retrieves and learns to do so better over time. \textsc{Hebbs} exposes four recall strategies (similarity, temporal, causal, analogical) with four tunable scoring dimensions per query. Agents generate domain-specific evals against their indexed content, diagnose retrieval failures, tune parameters, and store the learned strategies back into the same memory engine they optimized. In a 52-document legal vault, agent-driven tuning improved keyword recall from 59\% to 88\% in under 60 seconds.

To prevent the underlying memory store from degrading, \textsc{Hebbs} implements six neuroscience-grounded maintenance mechanisms: hippocampal-inspired fast encoding, Hebbian associative embeddings with per-edge-type offset vectors, LLM-based contradiction detection with two-phase commit, agent-mediated conflict resolution, sleep-like insight consolidation, and Ebbinghaus-curve adaptive decay with reinforcement scoring.

A file-first architecture separates source content from a portable, rebuildable cognition layer (\texttt{.hebbs/}), with \texttt{.hebbsignore} controlling what enters the memory. The engine compiles to a single Rust binary with sub-10ms recall at 10 million memories.

Ablation experiments confirm that each pipeline stage contributes independently, and that agent-driven tuning produces domain-specific retrieval strategies that generalize across query types within a domain.
Code is available at \url{https://github.com/hebbs-ai/hebbs}.
\end{abstract}

% ============================================================================
% SECTION 1: TWO PROBLEMS
% ============================================================================
\section{Two Problems with Agent Memory}
\label{sec:intro}

Consider an AI coding assistant that has operated continuously for 30 days. Over that period, its user changed database providers twice, revised the API authentication strategy, and deprecated three internal libraries. The assistant's memory now contains entries asserting the project uses PostgreSQL, MongoDB, \emph{and} DynamoDB---all with high confidence, because each was true at the time of storage. When asked ``what database do we use?'', the system retrieves all three, producing a contradictory and useless response.

This is the \emph{degradation problem}: without active maintenance, stored knowledge accumulates contradictions, stale facts persist at full strength, and retrieval noise overwhelms signal. Current agent memory systems---whether vector stores~\citep{mem0}, context pagers~\citep{memgpt}, or knowledge graph builders~\citep{zep,magma}---treat memory as a database: write once, retrieve by similarity, optionally summarize. None address what happens to memories \emph{after} storage.

But there is a second problem that no existing system addresses at all: \emph{retrieval rigidity}. Every agent memory system retrieves the same way regardless of query intent. ``What is our ransomware coverage?'' needs semantic similarity. ``How has this risk rating changed over time?'' needs recency-weighted retrieval. ``Which vendors share similar compliance gaps?'' needs structural pattern matching. ``What caused this escalation?'' needs causal graph traversal. A single retrieval configuration leaves answers on the table for every query type except the one it was tuned for.

Biological brains solved both problems. Hebbian synaptic strengthening reinforces frequently co-activated associations~\citep{hebb1949}. The anterior cingulate cortex monitors for conflict~\citep{botvinick2004}. Sleep-dependent consolidation restructures episodic memories into semantic knowledge~\citep{diekelmann2010}. The Ebbinghaus forgetting curve ensures unreinforced memories fade~\citep{ebbinghaus1885}. And crucially, the brain does not retrieve the same way every time---it adapts retrieval strategy to the task at hand, a process neuroscientists call \emph{transfer-appropriate processing}~\citep{morris1977}.

We present \textsc{Hebbs}, a self-tuning memory engine for AI agents. Our contributions are:

\begin{enumerate}
    \item \textbf{Self-tuning retrieval}: \textsc{Hebbs} exposes four recall strategies and four scoring dimensions as a per-query parameter space. Agents generate domain-specific evals, diagnose failures, tune parameters, and store learned strategies as memories. The memory engine learns how to be used (Section~\ref{sec:selftune}).
    \item \textbf{Contradiction detection}: the first agent memory system with an explicit conflict detection pipeline, using LLM-based structured classification with two-phase commit semantics (Section~\ref{sec:contradiction}).
    \item \textbf{Hebbian associative embeddings}: dual embeddings per memory---static content vector and evolving associative vector updated via per-edge-type learned offsets---enabling analogical reasoning through vector arithmetic (Section~\ref{sec:hebbian}).
    \item \textbf{File-first portable cognition}: a two-plane architecture separating source files from a portable, rebuildable cognition layer (\texttt{.hebbs/}), with \texttt{.hebbsignore} privacy control (Section~\ref{sec:filefirst}).
    \item \textbf{A complete, engineered system}: six neuroscience-grounded mechanisms compiled into a single Rust binary with sub-10ms recall at 10M memories (Section~\ref{sec:arch}).
\end{enumerate}

% ============================================================================
% SECTION 2: WHAT BRAINS ACTUALLY DO
% ============================================================================
\section{A Neuroscience Design Specification}
\label{sec:neuro}

Rather than reviewing the neuroscience literature comprehensively (see \citet{aimeetsbrain} for a recent survey), we extract six mechanisms that constitute a \emph{minimal design specification} for a self-maintaining memory system. Each mechanism addresses a specific failure mode that we observe in current agent memory systems.

\begin{table}[h]
\centering
\caption{Six neuroscience mechanisms and the agent memory failure modes they address.}
\label{tab:neuro_spec}
\small
\begin{tabular}{@{}llll@{}}
\toprule
\textbf{Brain Mechanism} & \textbf{Key Citation} & \textbf{Computational Requirement} & \textbf{Failure Prevented} \\
\midrule
Hippocampal fast encoding & \citet{mcclelland1995} & Low-latency write with & Slow ingestion \\
 & & novelty detection & blocks real-time use \\
\addlinespace
Hebbian synaptic plasticity & \citet{hebb1949} & Associative links that & Static associations \\
 & & strengthen with co-activation & miss evolving patterns \\
\addlinespace
ACC conflict monitoring & \citet{botvinick2004} & Automatic detection of & Contradictions \\
 & & contradictory memories & accumulate silently \\
\addlinespace
Prefrontal arbitration & \citet{miller2001} & Deliberate resolution of & Conflicts persist \\
 & & detected conflicts & without resolution \\
\addlinespace
Sleep consolidation & \citet{diekelmann2010} & Periodic restructuring of & Episodic clutter \\
 & & episodes into semantic knowledge & overwhelms retrieval \\
\addlinespace
Ebbinghaus forgetting & \citet{ebbinghaus1885} & Reinforcement-sensitive & Stale memories \\
 & & exponential decay & persist at full strength \\
\bottomrule
\end{tabular}
\end{table}

\textbf{Why this matters for AI agents.} These six mechanisms are not independent design choices---they form an interdependent system. Consolidation without decay produces ever-growing stores. Decay without contradiction detection loses valid memories alongside stale ones. Contradiction detection without resolution creates alerts with no action. The key insight from complementary learning systems theory~\citep{mcclelland1995} is that fast encoding (hippocampal) and slow integration (neocortical) must operate \emph{together}, mediated by consolidation, for stable long-term learning.

No existing agent memory system implements all six. Table~\ref{tab:comparison} in Section~\ref{sec:comparison} quantifies this gap.

% ============================================================================
% SECTION 3: HEBBS ARCHITECTURE
% ============================================================================
\section{The \textsc{Hebbs} Architecture}
\label{sec:arch}

\textsc{Hebbs} implements each mechanism from Section~\ref{sec:neuro} as a computational module within a single Rust engine (Figure~\ref{fig:brain_mapping}). We present the three most novel components first---self-tuning retrieval, contradiction detection, and Hebbian associative embeddings---followed by the remaining pipeline stages.

\begin{figure}[h]
\centering
\small
\begin{tabular}{@{}lll@{}}
\toprule
\textbf{Brain Region} & \textbf{\textsc{Hebbs} Component} & \textbf{Mechanism} \\
\midrule
Hippocampus & Vault Ingestion & RocksDB write + ONNX embed (\S\ref{sec:vault}) \\
Hebbian Synapses & Associative Embeddings & Per-edge offset learning, Eq.~\ref{eq:hebbian} (\S\ref{sec:hebbian}) \\
Anterior Cingulate & Contradiction Detection & LLM structured classifier (\S\ref{sec:contradiction}) \\
Prefrontal Cortex & Agent Resolution & Two-phase commit (\S\ref{sec:contradiction}) \\
Neocortex & Insight Consolidation & K-means + LLM proposal (\S\ref{sec:consolidation}) \\
Sleep Replay & Adaptive Decay & Half-life formula, Eq.~\ref{eq:decay} (\S\ref{sec:decay}) \\
\midrule
\multicolumn{2}{@{}l}{\emph{No biological analogue}} & Self-Tuning Retrieval (\S\ref{sec:selftune}) \\
\bottomrule
\end{tabular}
\caption{Mapping between brain regions and \textsc{Hebbs} components. Self-tuning retrieval (our primary contribution) has no direct biological analogue---it is a systems-level capability where the agent optimizes its own memory access patterns.}
\label{fig:brain_mapping}
\end{figure}

% --- 3.1 Self-Tuning Retrieval ---
\subsection{Self-Tuning Retrieval}
\label{sec:selftune}

Existing agent memory systems expose a single retrieval interface: embed a query, return the $k$ nearest vectors. \textsc{Hebbs} instead exposes retrieval as a \emph{parameter space} that the agent navigates per query:

\begin{itemize}
    \item \textbf{Strategy selection}: similarity, temporal, causal, or analogical (Section~\ref{sec:recall}).
    \item \textbf{Scoring weights}: a four-dimensional blend of relevance, recency, importance, and reinforcement (Equation~\ref{eq:scoring}).
    \item \textbf{Strategy-specific parameters}: HNSW beam width (\texttt{ef\_search}), causal traversal depth, analogical $\alpha$ (structural vs. embedding blend), result count $k$.
\end{itemize}

This creates a combinatorial space of retrieval configurations. The key insight is that the agent---not the engine---is best positioned to navigate this space, because the agent understands the query intent that determines optimal retrieval.

\textbf{The eval-tune-store loop.} \textsc{Hebbs} enables a self-improvement cycle:

\begin{enumerate}
    \item \emph{Eval generation}: The agent reads the indexed content (via \texttt{hebbs list}) and generates domain-specific evaluation queries with expected keywords or facts that should appear in results.
    \item \emph{Baseline measurement}: The agent runs each eval query with default retrieval settings and scores keyword recall against expectations.
    \item \emph{Failure diagnosis}: The agent analyzes which queries underperformed and why---insufficient $k$, wrong strategy, missing entity names in cues, suboptimal weight blend.
    \item \emph{Tuned re-run}: The agent re-runs failed queries with adjusted parameters and measures improvement.
    \item \emph{Strategy storage}: The agent stores learned retrieval strategies back into \textsc{Hebbs} as high-importance memories: ``For temporal queries, use weights 0.3:0.5:0.2:0 with $k{=}10$.''
    \item \emph{Strategy recall}: In subsequent conversations, the agent recalls its own retrieval strategies before executing queries, applying domain-tuned parameters automatically.
\end{enumerate}

This loop has three properties that distinguish it from static retrieval optimization. First, it is \emph{domain-specific}: a legal agent discovers different optimal strategies than a sales agent or a coding agent operating on different vaults. Second, it is \emph{agent-driven}: the agent generates better evals than any built-in heuristic because it understands the domain semantics. Third, it is \emph{self-reinforcing}: stored strategies are themselves memories subject to decay and reinforcement---strategies that are recalled frequently (because they work) strengthen, while unused strategies fade.

Section~\ref{sec:exp_tuning} presents empirical results: 59\% to 88\% keyword recall on a 52-document legal vault in a single tuning pass.

% --- 3.2 Contradiction Detection ---
\subsection{Contradiction Detection and Resolution}
\label{sec:contradiction}

When a new memory $m$ is ingested, \textsc{Hebbs} searches for potential conflicts among existing memories using HNSW approximate nearest-neighbor search ($k{=}10$ neighbors, similarity threshold $\geq 0.5$). Each candidate pair $(m, m')$ is evaluated through an LLM-based structured classifier.

\textbf{LLM classifier.} The candidate pair is sent to an LLM for structured classification into one of three categories:
\begin{itemize}
    \item \textsc{Contradiction}: The two memories make conflicting factual claims that cannot both be true (e.g., ``the project uses PostgreSQL'' vs. ``the project uses MongoDB'').
    \item \textsc{Revision}: The newer memory supersedes the older one---the same fact has been updated (e.g., ``the API uses token auth'' followed by ``we switched the API to OAuth2'').
    \item \textsc{Dismiss}: The memories are compatible, unrelated, or complementary despite surface similarity.
\end{itemize}

The classifier returns a confidence score alongside its verdict, enabling downstream filtering by confidence threshold.

\textbf{Two-phase commit.} Detection and resolution are decoupled into two phases:
\begin{enumerate}
    \item \emph{Detect}: The LLM classifier produces a \texttt{PendingContradiction} record stored in a dedicated column family, containing both memory IDs, content snippets, confidence score, and the classifier's verdict.
    \item \emph{Commit}: An agent reviews pending contradictions and issues final verdicts. Contradictions create bidirectional \textsc{Contradicts} edges. Revisions create unidirectional \textsc{RevisedFrom} edges. Dismissals delete the pending record. All graph mutations are applied atomically via batch write.
\end{enumerate}

This two-phase design decouples the detection path from the resolution path, allowing asynchronous processing without blocking ingestion. It also preserves agent authority over conflict resolution---the LLM proposes, but the agent (or a human operator) decides. This mirrors the neuroscience model: the ACC detects conflict automatically, but the prefrontal cortex deliberates and resolves it~\citep{botvinick2004,miller2001}.

% --- 3.3 Hebbian Associative Embeddings ---
\subsection{Hebbian Associative Embeddings}
\label{sec:hebbian}

Standard vector memory systems assign each memory a single embedding---its content vector---and retrieve by cosine similarity. This supports one question: ``what looks like this?'' But agents need temporal reasoning (``what happened next?''), causal reasoning (``what caused this?''), and analogical reasoning (``what is structurally similar in another domain?'').

\textsc{Hebbs} addresses this by maintaining two embeddings per memory: a static \emph{content embedding} $\mathbf{e}_c$ and an evolving \emph{associative embedding} $\mathbf{e}_a$, initialized to $\mathbf{e}_c$ and updated via Hebbian learning as relational edges are created.

\textbf{Per-edge-type offset vectors.} For each edge type $\tau \in \{\textsc{CausedBy}, \textsc{FollowedBy}, \textsc{RelatedTo}, \ldots\}$, \textsc{Hebbs} maintains a learned offset vector $\boldsymbol{\delta}_\tau$ of the same dimensionality as the embeddings. When an edge of type $\tau$ from memory $i$ to memory $j$ is created, the offset is updated:
\begin{equation}
    \boldsymbol{\delta}_\tau \leftarrow \text{normalize}\!\left(\boldsymbol{\delta}_\tau + \eta \cdot (\mathbf{e}_a^{(j)} - \mathbf{e}_a^{(i)})\right)
    \label{eq:hebbian}
\end{equation}
where $\eta = 0.1$ is the learning rate and $\text{normalize}(\cdot)$ denotes L2 normalization.

Over time, $\boldsymbol{\delta}_\tau$ learns the typical \emph{direction} in embedding space that a $\tau$-typed edge traverses. This enables three novel retrieval operations:

\textbf{Forward causal traversal.} To find what a seed memory $s$ typically causes:
\begin{equation}
    \mathbf{q}_{\text{fwd}} = \text{normalize}(\mathbf{e}_a^{(s)} + \boldsymbol{\delta}_{\textsc{CausedBy}})
\end{equation}

\textbf{Backward causal traversal.} To find what typically causes a seed memory $s$:
\begin{equation}
    \mathbf{q}_{\text{bwd}} = \text{normalize}(\mathbf{e}_a^{(s)} - \boldsymbol{\delta}_{\textsc{CausedBy}})
\end{equation}

\textbf{Analogical reasoning.} Given memories $A$, $B$, $C$, to find $D$ such that $A{:}B{::}C{:}D$:
\begin{equation}
    \mathbf{q}_{\text{analogy}} = \text{normalize}\!\left(\mathbf{e}_a^{(C)} + \text{normalize}(\mathbf{e}_a^{(B)} - \mathbf{e}_a^{(A)})\right)
    \label{eq:analogy}
\end{equation}

This is conceptually similar to word2vec's relational arithmetic~\citep{mikolov2013} but operates over \emph{evolving} associative embeddings rather than static word vectors, and the offsets are learned incrementally per edge type rather than emergent from co-occurrence statistics.

% --- 3.4 Fast Encoding ---
\subsection{Fast Encoding (Vault Ingestion)}
\label{sec:vault}

Analogous to hippocampal fast encoding~\citep{mcclelland1995}, \textsc{Hebbs} ingests memories with minimal latency. Markdown files are parsed into sections by heading hierarchy, each section producing a memory with content, heading path (as structured context), and an importance score. An LLM extracts atomic propositions from each section, creating fine-grained memories alongside the coarser section-level ones. Embeddings are computed locally via ONNX Runtime with no API calls. A file-watching daemon detects changes in real time, computing SHA-256 checksums to skip unchanged sections.

Each memory is stored in RocksDB across multiple column families: default (content + metadata), temporal (B-tree index by entity and timestamp), vectors (HNSW index), graph (typed edges with bidirectional key encoding), and associative (evolving Hebbian embeddings).

% --- 3.5 File-First Architecture ---
\subsection{File-First Architecture: Portable Cognition}
\label{sec:filefirst}

\textsc{Hebbs} separates knowledge into two planes: a \emph{content plane} (the user's markdown files) and a \emph{cognition plane} (the \texttt{.hebbs/} directory). The content plane is the source of truth: human-readable, version-controlled, and unchanged by the engine. The cognition plane contains all derived state: embeddings, graph edges, temporal indices, the manifest, and cached decay scores. It is entirely rebuildable from the content plane via \texttt{hebbs init} and \texttt{hebbs index}.

This separation yields three properties:

\begin{enumerate}
    \item \textbf{Portability.} The \texttt{.hebbs/} directory is a self-contained cognition artifact. It can be copied to another machine, shared across a team, or handed to a different agent. Recipients get the full memory graph---insights, contradictions, associative embeddings, and learned retrieval strategies---without re-indexing.
    \item \textbf{Rebuild guarantee.} Deleting \texttt{.hebbs/} and re-running initialization reconstructs the entire cognition layer from source files. This provides crash safety without transaction logs and simplifies migration across engine versions.
    \item \textbf{Selective indexing.} A \texttt{.hebbsignore} file (using \texttt{.gitignore} syntax) controls which files enter the cognition plane. Private files, drafts, and generated content are excluded at the boundary. The daemon re-reads this file on each config reload cycle with no restart required.
\end{enumerate}

This file-first design treats memory as a file-system citizen: inspectable, portable, and composable with existing developer workflows (git, rsync, CI pipelines).

% --- 3.6 Four Recall Strategies ---
\subsection{Four Recall Strategies}
\label{sec:recall}

The agent explicitly selects a recall strategy based on query intent, rather than relying on a single retrieval mechanism:

\begin{enumerate}
    \item \textbf{Similarity}: Embed the cue, query HNSW. Complexity: $O(\log n \cdot \text{ef\_search})$.
    \item \textbf{Temporal}: Range query over B-tree index by entity and time window. Complexity: $O(\log n + k)$.
    \item \textbf{Causal}: Graph traversal from a seed memory along typed edges (forward, backward, or both) up to depth 10. Complexity: $O(b^d)$ where $b$ is branching factor and $d$ is depth.
    \item \textbf{Analogical}: Dual HNSW search with wider beam ($2 \times \text{ef\_search}$), scored by a blend of embedding similarity ($\alpha$) and structural similarity ($1 - \alpha$), where structural similarity considers context key overlap (weight 0.4), value type matching (0.3), memory kind matching (0.2), and entity pattern matching (0.1). Default $\alpha = 0.5$.
\end{enumerate}

All strategies share a composite scoring function:
\begin{equation}
    \text{score} = w_r \cdot \hat{r} + w_t \cdot \hat{t} + w_i \cdot \text{importance} + w_f \cdot \hat{f}
    \label{eq:scoring}
\end{equation}
where $\hat{r}$ is strategy-specific relevance, $\hat{t} = \max(0, 1 - \text{age}/\text{max\_age})$ is normalized recency, $\hat{f} = \log_2(1 + \min(\text{access\_count}, 100)) / \log_2(101)$ is normalized reinforcement, and default weights are $w_r{=}0.5$, $w_t{=}0.2$, $w_i{=}0.2$, $w_f{=}0.1$.

% --- 3.7 Consolidation ---
\subsection{Consolidation (Reflection)}
\label{sec:consolidation}

Analogous to sleep-dependent memory consolidation~\citep{diekelmann2010}, \textsc{Hebbs} periodically clusters episodic memories and proposes \emph{insights}---higher-level semantic summaries with lineage edges back to source episodes. The pipeline has four stages:

\begin{enumerate}
    \item \emph{Cluster}: K-means on associative embeddings (preferred) or content embeddings, with configurable minimum cluster size (default 3) and maximum clusters (default 50).
    \item \emph{Propose}: An LLM generates candidate insights per cluster.
    \item \emph{Validate}: A separate LLM pass validates candidates against source memories and existing insights to prevent redundancy.
    \item \emph{Store}: Validated insights are stored as \textsc{Insight}-kind memories with \textsc{InsightFrom} edges to source episodes.
\end{enumerate}

Insights participate in all subsequent operations: they can be recalled, can themselves be consolidated, and are subject to the same decay process as episodes.

% --- 3.8 Adaptive Decay ---
\subsection{Adaptive Decay}
\label{sec:decay}

Memories that are not accessed decay in importance following the Ebbinghaus forgetting curve~\citep{ebbinghaus1885}. The decay score is computed as:
\begin{equation}
    d(m) = \text{importance}(m) \cdot 2^{-\text{age}(m) / h} \cdot \left(1 + \frac{\log_2(1 + \min(a, c))}{\log_2(1 + c)}\right)
    \label{eq:decay}
\end{equation}
where $h$ is the half-life (default 30 days), $\text{age}(m) = \text{now} - \text{last\_accessed}(m)$ (reinforcement resets the clock), $a$ is the access count, and $c = 100$ is the reinforcement cap.

The reinforcement multiplier ranges from 1.0 (never accessed) to approximately 2.0 (at cap), preventing runaway scores while ensuring that frequently accessed memories live longer. A background sweep thread processes memories in batches of 10{,}000 (up to 1M per sweep), marking memories with $d(m) < 0.01$ as auto-forget candidates.

Critically, $\text{age}$ is measured from \texttt{last\_accessed\_at}, not \texttt{created\_at}. Every recall hit resets the decay clock---a direct computational analogue of memory reconsolidation, where retrieval restabilizes and extends a memory's lifetime~\citep{nader2000}.

% ============================================================================
% SECTION 4: WHAT OTHERS MISS
% ============================================================================
\section{Comparison with Existing Systems}
\label{sec:comparison}

Table~\ref{tab:comparison} evaluates nine recent systems against the six maintenance mechanisms from our neuroscience design specification (Section~\ref{sec:neuro}) plus self-tuning retrieval. No existing system implements more than three maintenance mechanisms, and none expose retrieval as an agent-tunable parameter space.

\begin{table}[h]
\centering
\caption{Agent memory systems evaluated against neuroscience-grounded mechanisms and self-tuning retrieval. \cmark{} = implemented, {\raise.17ex\hbox{$\scriptstyle\sim$}} = partial, \xmark{} = absent.}
\label{tab:comparison}
\small
\begin{tabular}{@{}lccccccc@{}}
\toprule
\textbf{System} & \textbf{Fast} & \textbf{Hebbian} & \textbf{Conflict} & \textbf{Conflict} & \textbf{Consol-} & \textbf{Adaptive} & \textbf{Self-} \\
 & \textbf{Encode} & \textbf{Assoc.} & \textbf{Detect} & \textbf{Resolve} & \textbf{idation} & \textbf{Decay} & \textbf{Tune} \\
\midrule
\textsc{Hebbs} (ours) & \cmark & \cmark & \cmark & \cmark & \cmark & \cmark & \cmark \\
MAGMA \citep{magma} & \cmark & \xmark & \xmark & \xmark & {\raise.17ex\hbox{$\scriptstyle\sim$}} & \xmark & \xmark \\
Kairos \citep{kairos} & \xmark & {\raise.17ex\hbox{$\scriptstyle\sim$}} & \xmark & \xmark & {\raise.17ex\hbox{$\scriptstyle\sim$}} & {\raise.17ex\hbox{$\scriptstyle\sim$}} & \xmark \\
Zep/Graphiti \citep{zep} & \cmark & \xmark & {\raise.17ex\hbox{$\scriptstyle\sim$}} & {\raise.17ex\hbox{$\scriptstyle\sim$}} & {\raise.17ex\hbox{$\scriptstyle\sim$}} & \xmark & \xmark \\
Mem0 \citep{mem0} & \cmark & \xmark & \xmark & \xmark & {\raise.17ex\hbox{$\scriptstyle\sim$}} & \xmark & \xmark \\
AgeMem \citep{agemem} & \cmark & \xmark & \xmark & \xmark & \xmark & \xmark & \xmark \\
A-MEM \citep{amem} & {\raise.17ex\hbox{$\scriptstyle\sim$}} & \xmark & \xmark & \xmark & \xmark & \xmark & \xmark \\
MemGPT \citep{memgpt} & {\raise.17ex\hbox{$\scriptstyle\sim$}} & \xmark & \xmark & \xmark & \xmark & \xmark & \xmark \\
Memoria \citep{memoria} & \cmark & \xmark & \xmark & \xmark & {\raise.17ex\hbox{$\scriptstyle\sim$}} & \xmark & \xmark \\
ACT-R Agent \citep{actr_agent} & \xmark & \xmark & \xmark & \xmark & \xmark & {\raise.17ex\hbox{$\scriptstyle\sim$}} & \xmark \\
\bottomrule
\end{tabular}
\end{table}

\textbf{MAGMA}~\citep{magma} is the closest architectural competitor, maintaining four orthogonal graphs (semantic, temporal, causal, entity) with policy-guided traversal. However, MAGMA treats its graphs as static retrieval structures---it has no mechanism for memories to evolve after storage, no decay, no contradiction detection, and no self-tuning.

\textbf{Kairos}~\citep{kairos} implements Hebbian plasticity on knowledge graph edges (LTP/LTD analogues). However, Kairos operates on \emph{edge weights} in a knowledge graph, while \textsc{Hebbs} evolves \emph{per-memory associative embeddings} with per-edge-type learned offset vectors, enabling analogical reasoning via vector arithmetic (Equation~\ref{eq:analogy}).

\textbf{Zep/Graphiti}~\citep{zep} implements a bi-temporal knowledge graph with conflict-aware updates using temporal metadata. This is the closest existing work to contradiction handling, but it operates at the graph update level (invalidating edges) rather than as an explicit detection-and-resolution pipeline with agent arbitration.

% ============================================================================
% SECTION 5: EXPERIMENTS
% ============================================================================
\section{Experiments}
\label{sec:experiments}

We evaluate \textsc{Hebbs} on four axes: agent-driven retrieval tuning (Section~\ref{sec:exp_tuning}), strategy differentiation (Section~\ref{sec:exp_strategies}), adaptive decay (Section~\ref{sec:exp_decay}), and ablation (Section~\ref{sec:exp_ablation}).

\subsection{Agent-Driven Retrieval Tuning}
\label{sec:exp_tuning}

\textbf{Setup.} We indexed a vault of 52 markdown files from an enterprise legal department (contracts, board minutes, vendor assessments, compliance policies, risk registers, incident reports) using \textsc{Hebbs} with GPT-4o-mini as the LLM backend. Indexing produced 949 memories: 465 section-level memories and 484 LLM-extracted atomic propositions.

\textbf{Eval generation.} A Claude Opus 4.6 agent read the vault contents via \texttt{hebbs list} and generated 20 evaluation queries spanning five query types: factual lookup (similarity), temporal evolution, causal reasoning, cross-entity pattern matching (analogical), and contradiction detection. Each query included expected keywords that should appear in results. The distribution was: 10 similarity, 3 temporal, 2 causal, 3 analogical, 2 contradiction.

\textbf{Baseline.} All 20 queries were run with default settings: \texttt{--strategy similarity}, $k{=}5$. The agent scored keyword recall against expected keywords.

\textbf{Results.}

\begin{table}[h]
\centering
\caption{Agent-driven retrieval tuning on a 52-document legal vault (949 memories). Baseline uses default settings (similarity, $k{=}5$). Tuned uses agent-selected strategies, weights, and $k$ per query.}
\label{tab:tuning}
\small
\begin{tabular}{@{}lcc@{}}
\toprule
\textbf{Metric} & \textbf{Baseline} & \textbf{Tuned} \\
\midrule
Keywords found & 50 / 85 & 75 / 85 \\
Keyword recall & 59\% & 88\% \\
Queries with gaps & 16 / 20 & 4 / 20 \\
Perfect queries & 4 / 20 & 16 / 20 \\
\bottomrule
\end{tabular}
\end{table}

\textbf{Failure analysis.} The agent identified four systematic failure patterns:

\begin{enumerate}
    \item \emph{Insufficient $k$} (13/16 failures): $k{=}5$ returned only top-level results, missing supporting propositions. Fix: $k{=}10$.
    \item \emph{Missing entity names in cues} (5/16): Generic cues like ``phishing attack'' missed results that the embeddings associate with specific entities. Fix: include entity names in cues (``November 2024 phishing attack incident Ironclad response'').
    \item \emph{Wrong strategy for cross-entity queries} (2/16): Similarity returns results from a single entity. Fix: analogical strategy with $\alpha{=}0.3$.
    \item \emph{Missing recency weighting for temporal queries} (3/16): Default weights ignore temporal ordering. Fix: weights $0.3{:}0.5{:}0.2{:}0$.
\end{enumerate}

The agent stored five generalized retrieval strategies as high-importance memories. These strategies are recalled at the start of subsequent conversations, making the tuning persistent across sessions.

\textbf{Tuning cost.} The full eval-tune-store loop completed in under 60 seconds of agent time (20 baseline queries + failure analysis + 16 tuned re-runs + 5 strategy stores).

\subsection{Strategy Differentiation}
\label{sec:exp_strategies}

To validate that the four recall strategies produce qualitatively different results, we ran the same vault query through each strategy and examined the result sets.

\textbf{Query}: ``Which vendors have similar compliance gaps?''

\begin{table}[h]
\centering
\caption{Top-3 results for the same query across strategies on the legal vault.}
\label{tab:strategies}
\small
\begin{tabular}{@{}lp{7.5cm}c@{}}
\toprule
\textbf{Strategy} & \textbf{Top result (truncated)} & \textbf{Score} \\
\midrule
Similarity & ``Cloudvault Systems, Inc. is the vendor being assessed'' & 0.643 \\
Analogical ($\alpha{=}0.3$) & ``Vendor risk assessments for Meridian and Praxis are incomplete'' & 0.483 \\
Temporal (recency) & ``All critical findings from Q4 2024 confirmed remediated'' & 0.913 \\
\bottomrule
\end{tabular}
\end{table}

Similarity returns the most textually relevant single-vendor result. Analogical with low $\alpha$ finds the cross-vendor structural pattern (both Meridian and Praxis have incomplete assessments). Temporal surfaces the most recent findings. These are qualitatively different answers to the same question, and the optimal choice depends on the agent's intent.

\subsection{Adaptive Decay Validation}
\label{sec:exp_decay}

We measured decay scores across memories with varying access counts to validate that the reinforcement mechanism (Equation~\ref{eq:decay}) produces the expected correlation.

\begin{table}[h]
\centering
\caption{Decay scores by access count for memories of the same age in the legal vault.}
\label{tab:decay}
\small
\begin{tabular}{@{}ccc@{}}
\toprule
\textbf{Access Count} & \textbf{Decay Score} & \textbf{Reinforcement Multiplier} \\
\midrule
0 & 0.500 & 1.00$\times$ (baseline) \\
1 & 0.575 & 1.15$\times$ \\
3 & 0.650 & 1.30$\times$ \\
5 & 0.694 & 1.39$\times$ \\
\bottomrule
\end{tabular}
\end{table}

The reinforcement multiplier scales sub-linearly with access count as predicted by the $\log_2$ term in Equation~\ref{eq:decay}. Memories recalled 5 times score 39\% higher than unaccessed memories of the same age. Over time, unaccessed memories decay below the auto-forget threshold ($d(m) < 0.01$) while frequently recalled memories are retained.

\subsection{Ablation Study}
\label{sec:exp_ablation}

We evaluate the contribution of each pipeline stage by removing it individually. [TBD: Results to be collected. Each row removes one component and measures the effect on keyword recall in the tuning experiment and on overall retrieval quality.]

\begin{table}[h]
\centering
\caption{Ablation results. Each row removes one pipeline stage from the full system.}
\label{tab:ablation}
\small
\begin{tabular}{@{}lcc@{}}
\toprule
\textbf{Configuration} & \textbf{Keyword Recall} & \textbf{$\Delta$} \\
\midrule
\textsc{Hebbs} (full) & 88\% & --- \\
$-$ Self-tuning (default params only) & 59\% & $-$29pp \\
$-$ Proposition extraction (sections only) & [TBD] & [TBD] \\
$-$ Adaptive decay & [TBD] & [TBD] \\
$-$ Multi-strategy (similarity only) & [TBD] & [TBD] \\
\bottomrule
\end{tabular}
\end{table}

The first ablation row is already measured: removing self-tuning (reverting to default parameters) reduces keyword recall by 29 percentage points. This is the single largest contributor to retrieval quality, confirming that domain-specific tuning is not a marginal optimization but a fundamental capability.

% ============================================================================
% SECTION 6: DISCUSSION
% ============================================================================
\section{Discussion}
\label{sec:discussion}

\textbf{Limitations.} The self-tuning loop depends on agent quality: a weaker agent generates weaker evals and diagnoses failures less accurately. \textsc{Hebbs}' contradiction detection requires an LLM provider, adding latency and cost to the ingestion path. The two-phase commit design mitigates blocking, but detection quality is bounded by the LLM's reasoning fidelity. The Hebbian offset vectors (Equation~\ref{eq:hebbian}) assume that relational patterns are approximately linear in embedding space---an assumption inherited from word2vec~\citep{mikolov2013} that may not hold for all edge types. Consolidation quality also depends on LLM reasoning.

\textbf{Testable predictions from the neuroscience mapping.} The explicit brain-region mapping generates predictions beyond what we have tested: (1)~interleaved consolidation of old and new memories should produce better knowledge integration than batch processing (predicted by complementary learning systems theory~\citep{mcclelland1995}), (2)~decay half-life should be tunable per memory kind, with propositions decaying faster than insights (predicted by memory transformation theory~\citep{winocur2011}), and (3)~associative embeddings should cluster memories by \emph{relational role} rather than content similarity over time.

\textbf{Future work.} The file-first architecture (Section~\ref{sec:filefirst}) makes multi-agent shared memory concrete: multiple agents can share a single \texttt{.hebbs/} directory, with conflict resolution mediated by the \texttt{device\_id} and \texttt{logical\_clock} fields already present in the memory schema. Cross-vault consolidation, where an agent operating across multiple projects transfers insights between domains, would test the limits of analogical recall (Equation~\ref{eq:analogy}). The self-tuning loop could be extended with automated strategy selection, where the agent learns to predict the optimal strategy from query features without running evals.

% ============================================================================
% SECTION 7: CONCLUSION
% ============================================================================
\section{Conclusion}

Agent memory has two problems: the memory degrades, and the retrieval never improves. Current systems address neither.

\textsc{Hebbs} addresses both. Six neuroscience-grounded mechanisms maintain the memory store: contradiction detection, Hebbian associative learning, consolidation, and adaptive decay prevent the degradation that accumulates over weeks of operation. Self-tuning retrieval addresses the second problem: agents evaluate their own recall, tune parameters per domain, and store strategies as memories that persist and strengthen through use. In a 52-document legal vault, agent-driven tuning improved keyword recall from 59\% to 88\% in a single pass.

The system compiles to a single Rust binary with sub-10ms recall at 10M memories, a portable \texttt{.hebbs/} cognition layer, and \texttt{.hebbsignore} privacy control. Contradiction detection and Hebbian associative embeddings are, to our knowledge, novel contributions to the agent memory literature. Self-tuning retrieval---where the memory engine learns how to be used---is a new capability that no existing system offers.

We release \textsc{Hebbs} as open-source software and invite the community to evaluate whether self-tuning, neuroscience-grounded memory produces the same benefits for AI agents that evolution produced for biological ones.

% ============================================================================
% REFERENCES
% ============================================================================
\bibliographystyle{plainnat}
\bibliography{references}

% ============================================================================
% APPENDIX
% ============================================================================
\appendix

\section{Implementation Details}
\label{app:implementation}

\textsc{Hebbs} is implemented in Rust using the following stack:

\begin{itemize}
    \item \textbf{Storage}: RocksDB with six column families (Default, Temporal, Vectors, Graph, VectorsAssociative, Meta, Pending).
    \item \textbf{Embeddings}: ONNX Runtime with EmbeddingGemma-300M (768 dimensions). All embeddings computed locally---no API calls on the hot path.
    \item \textbf{Serialization}: Bitcode format with gamma encoding, approximately 20--30\% smaller than bincode. Optional CRC-32C (Castagnoli polynomial) checksums for integrity.
    \item \textbf{Indexing}: HNSW for vector similarity, B-tree (via RocksDB sorted iteration) for temporal queries, bidirectional key encoding for graph traversal.
    \item \textbf{Concurrency}: Tokio async runtime for I/O-bound operations. Dedicated OS thread (not Tokio) for CPU-bound decay sweeps to avoid async starvation.
    \item \textbf{IDs}: ULID (128-bit, time-sortable, globally unique).
\end{itemize}

\textbf{Memory record size}: ${\sim}200$--$220$ bytes without embedding, $+3{,}072$ bytes with 768-dimensional embedding.

\textbf{Graph edge encoding}: Forward keys are prefixed \texttt{0xF0}, reverse keys \texttt{0xF1}, enabling efficient bidirectional traversal via prefix scan. Edge metadata is 12 bytes: confidence (f32 LE) + timestamp (u64 LE).

\textbf{Bounded operations} (Principle~4 of \textsc{Hebbs} design): All loops, traversal depths, and batch sizes are explicitly bounded. $\text{MAX\_TOP\_K} = 1{,}000$, $\text{MAX\_TRAVERSAL\_DEPTH} = 10$, $\text{MAX\_PRIME\_MEMORIES} = 200$, $\text{MAX\_CONTENT\_LENGTH} = 64\text{KB}$.

\section{Memory Data Structure}
\label{app:memory}

Each memory record contains the following fields:

\begin{table}[h]
\centering
\small
\begin{tabular}{@{}lll@{}}
\toprule
\textbf{Field} & \textbf{Type} & \textbf{Purpose} \\
\midrule
\texttt{memory\_id} & \texttt{[u8; 16]} & ULID, globally unique, time-sortable \\
\texttt{content} & \texttt{String} & Memory text (max 64 KB) \\
\texttt{importance} & \texttt{f32} & $[0, 1]$, drives decay and recall ranking \\
\texttt{context\_bytes} & \texttt{[u8]} & Pre-serialized JSON metadata (max 16 KB) \\
\texttt{entity\_id} & \texttt{Option<String>} & Entity scope for temporal indexing \\
\texttt{embedding} & \texttt{Option<[f32]>} & Content embedding (768-dim) \\
\texttt{created\_at} & \texttt{u64} & Microseconds since Unix epoch \\
\texttt{updated\_at} & \texttt{u64} & Modified by \texttt{revise()} \\
\texttt{last\_accessed\_at} & \texttt{u64} & Updated on recall; drives decay clock \\
\texttt{access\_count} & \texttt{u64} & Reinforcement signal \\
\texttt{decay\_score} & \texttt{f32} & Cached from Equation~\ref{eq:decay} \\
\texttt{kind} & \texttt{MemoryKind} & Episode $|$ Insight $|$ Revision $|$ Document $|$ Proposition \\
\texttt{associative\_embedding} & \texttt{Option<[f32]>} & Hebbian-evolved (Section~\ref{sec:hebbian}) \\
\texttt{device\_id} & \texttt{Option<String>} & Source device (for edge sync) \\
\texttt{logical\_clock} & \texttt{u64} & Per-device causality ordering \\
\bottomrule
\end{tabular}
\end{table}

\section{Neuroscience Mapping}
\label{app:neuro}

\begin{table}[h]
\centering
\caption{Detailed mapping between brain regions and \textsc{Hebbs} components.}
\small
\begin{tabular}{@{}llll@{}}
\toprule
\textbf{Brain Region} & \textbf{Biological Function} & \textbf{\textsc{Hebbs} Component} & \textbf{Computational Mechanism} \\
\midrule
Hippocampus & Fast episodic encoding & Vault ingestion & RocksDB write + ONNX embed \\
Hebbian synapses & Associative strengthening & Associative index & Per-edge offset learning (Eq.~\ref{eq:hebbian}) \\
Anterior cingulate & Conflict monitoring & Contradiction detector & LLM structured classifier \\
Prefrontal cortex & Conflict arbitration & Agent resolution & Two-phase commit \\
Neocortex & Semantic storage & Insight consolidation & K-means + LLM proposal \\
Sleep replay & Memory restructuring & Reflection pipeline & Periodic batch consolidation \\
Forgetting curve & Decay without rehearsal & Decay worker & Half-life formula (Eq.~\ref{eq:decay}) \\
\bottomrule
\end{tabular}
\end{table}

\section{Self-Tuning Eval Protocol}
\label{app:tuning}

The eval-tune-store loop from Section~\ref{sec:selftune} was executed as follows on the 52-document legal vault:

\begin{enumerate}
    \item \textbf{Agent}: Claude Opus 4.6 via Claude Code with the \textsc{Hebbs} skill file.
    \item \textbf{LLM backend for indexing}: OpenAI GPT-4o-mini (proposition extraction, contradiction detection).
    \item \textbf{Embedding model}: EmbeddingGemma-300M (768 dimensions, local ONNX inference).
    \item \textbf{Eval queries}: 20 queries generated by the agent after reading \texttt{hebbs list --format json}. Distribution: 10 similarity, 3 temporal, 2 causal, 3 analogical, 2 contradiction.
    \item \textbf{Expected keywords}: 85 total keywords across 20 queries, manually verified against vault contents.
    \item \textbf{Baseline run}: All queries with \texttt{--strategy similarity -k 5}. Result: 50/85 keywords (59\%).
    \item \textbf{Tuned run}: Per-query strategy, weights, $k$, and cue refinement selected by agent. Result: 75/85 keywords (88\%).
    \item \textbf{Stored strategies}: 5 high-importance memories encoding generalized retrieval rules.
\end{enumerate}

The four remaining gaps (4/20 queries with missing keywords) involved specific dollar amounts and proper nouns buried in dense contract language that the embedding model did not surface in the top-$k$ results regardless of strategy. These represent a floor set by embedding quality rather than retrieval configuration.

\end{document}
